We separate the causal target from the observed-law calculations throughout.
The target is \(Q_M(x,y)=P_M(i_Y=y\mid do(i_X=x))\); the quantities
\(e_z\), \(m_{zx}\), \(L_x\), and \(M_x\) are functions of the observed
law only.  We prove the asserted relation between them.

First note that every displayed conditional is well defined.  If
\(p(z)>0\), then nonnegativity and marginalization give
\[
 e_z(x)\geq0,
 \qquad
 \sum_{x\in\mathcal X}e_z(x)
 =\frac{\sum_xp(z,x)}{p(z)}=1.
\tag{1}
\]
If \(p(z,x)>0\), then necessarily \(p(z)>0\), and similarly
\[
 m_{zx}(y)\geq0,
 \qquad
 \sum_{y\in\mathcal Y}m_{zx}(y)
 =\frac{\sum_yp(z,x,y)}{p(z,x)}=1.
\tag{2}
\]
Thus \(e_z\in\Delta(\mathcal X)\), and every reached
\(m_{zx}\in\Delta(\mathcal Y)\).

We next prove necessity of the profile-count condition.  Let
\(M\in\mathfrak M(G,\mathcal A,V,W,T)\) induce \(p\).  By
qualitative functionality there is a map
\(f:\mathcal Z\to\mathcal B\) such that
\(\theta_{i_A}(a\mid z)={\bf 1}\{a=f(z)\}\).  Introduce the row notation
\[
 \pi(z)=\theta_{i_Z}(z),\qquad
 r_a(x)=\theta_{i_X}(x\mid a),\qquad
 g_{a,x}(y)=\theta_{i_Y}(y\mid a,x).
\]
The recursive causal factorization and summation over the hidden coordinate
give, cell by cell,
\[
 p(z,x,y)=\pi(z)r_{f(z)}(x)g_{f(z),x}(y).
\tag{3}
\]
Summing (3) first over \(y\), and then over \(x,y\), yields
\[
 p(z)=\pi(z),
 \qquad
 p(z,x)=p(z)r_{f(z)}(x).
\tag{4}
\]
For \(p(z)>0\), division in (4) is licensed and gives
\(e_z=r_{f(z)}\).  For \(p(z,x)>0\), (4) implies
\(r_{f(z)}(x)>0\), so division in (3) is licensed and gives
\[
 m_{zx}=g_{f(z),x}.
\tag{5}
\]
Consequently, if two positive-mass values \(z,z'\) satisfy
\(f(z)=f(z')\), then their entire \(e\)-vectors agree; hence their reached
supports agree, and (5) shows that their reached \(m\)-vectors agree.
Therefore \(\rho_z=\rho_{z'}\).  The map \(f\) can distinguish at most
\(\lvert\mathcal B\rvert\) profiles, proving
\[
 \lvert\mathcal R(p)\rvert\leq\lvert\mathcal B\rvert.
\tag{6}
\]

For sufficiency, assume (6) and choose an injection
\(\kappa:\mathcal R(p)\to\mathcal B\).  For every positive-mass \(z\),
set \(f(z)=\kappa(\rho_z)\); assign zero-mass \(z\)'s arbitrarily in
\(\mathcal B\).  Set \(\theta_{i_Z}(z)=p(z)\).  If
\(a=\kappa(\rho)\) is assigned, define
\[
 \theta_{i_X}(\cdot\mid a)=e_\rho,
 \qquad
 \theta_{i_Y}(\cdot\mid a,x)=m_{\rho x}
 \quad\text{when }e_\rho(x)>0,
\tag{7}
\]
where the profile equality makes both right-hand sides independent of the
representative \(z\).  Fill every row not fixed by (7) with an arbitrary
probability vector.  Make the \(A\)-row at \(z\) the point mass at
\(f(z)\).  This is a recursive causal Bayesian-network model and the
\(A\)-mechanism obeys qualitative functionality.  Its observed cells equal
\(p\): if \(p(z)=0\), both sides are zero; if \(p(z)>0\) but
\(p(z,x)=0\), then \(e_z(x)=0\), so both sides are zero; and if
\(p(z,x)>0\), equations (7), (1), and (2) give
\[
 p(z)e_z(x)m_{zx}(y)
 =p(z)\frac{p(z,x)}{p(z)}
       \frac{p(z,x,y)}{p(z,x)}
 =p(z,x,y).
\tag{8}
\]
All divisions in (8) are guarded by the positive denominators stated in
their cases.  The constructed marginal of \(X\) is the prescribed one,
and hence is positive at every \(x\) by hypothesis.  Thus joint-treatment
positivity for the singleton treatment is satisfied.  The causal
factorization and truncated-intervention semantics hold by the construction
of the recursive model, so Definition~\(\mathrm{def:model\text{-}class}\)
places it in \(\mathfrak M(G,\mathcal A,V,W,T)\).  This proves sufficiency.
If \(p\) is rational, all guarded ratios in (7) are rational; choosing, for
example, point masses for arbitrary rows makes every parameter rational.

We now characterize the causal fiber.  Fix any compatible model.  Under
intervention on \(i_X\), the truncated-intervention assumption removes the
factor \(r_{f(z)}(x)\), while the root and functional mechanisms remain.
Therefore
\[
 Q_M(x,y)=\sum_{z\in\mathcal Z}p(z)g_{f(z),x}(y).
\tag{9}
\]
For reached pairs, equation (5) identifies the outcome row.  Let
\(H_x=\{z:p(z)>0,\ p(z,x)=0\}\).  Splitting (9) into reached pairs,
holes, and zero-mass root values gives
\[
 Q_M(x,\cdot)
 =L_x+\sum_{z\in H_x}p(z)g_{f(z),x}.
\tag{10}
\]
Every summand row in the last term belongs to
\(\Delta(\mathcal Y)\), and its weights sum to \(M_x\).  Hence, when
\(M_x>0\), the normalized sum belongs to \(\Delta(\mathcal Y)\); when
\(M_x=0\), the sum is zero.  Equation (10) proves the outer inclusion
\[
 \mathcal Q(p)\subseteq
 \prod_{x\in\mathcal X}\bigl(L_x+M_x\Delta(\mathcal Y)\bigr).
\tag{11}
\]
Notice also, using (2), that
\[
 \sum_{y\in\mathcal Y}L_x(y)
 =\sum_{z:p(z,x)>0}p(z)=1-M_x.
\tag{12}
\]

For the reverse inclusion, choose an arbitrary vector in the right-hand
side of (11).  Thus, for each \(x\), write
\(q_x=L_x+M_xs_x\) with \(s_x\in\Delta(\mathcal Y)\).  Equivalently,
when \(M_x>0\), the mechanically computable choice is
\(s_x=(q_x-L_x)/M_x\); the membership assumption and (12) prove that its
coordinates are nonnegative and sum to one.  When \(M_x=0\), choose any
\(s_x\in\Delta(\mathcal Y)\).  In the sufficiency construction (7), set
\[
 \theta_{i_Y}(\cdot\mid\kappa(\rho),x)=s_x
 \quad\text{for every assigned profile }\rho\text{ with }e_\rho(x)=0.
\tag{13}
\]
This is well defined because all members of a profile have the same
\(e\)-vector.  It does not change the observed law: every row altered in
(13) is multiplied observationally by \(e_\rho(x)=0\).  Equations
(10) and (13) give
\[
 Q_M(x,\cdot)=L_x+\sum_{z\in H_x}p(z)s_x
 =L_x+M_xs_x=q_x.
\tag{14}
\]
The rows in (13) have distinct indices \((a,x)\) as \(x\) varies, so all
treatment-specific choices are simultaneous.  This proves the reverse
inclusion, sharpness, and original-alphabet attainment.  If \(p\) and
\(q\) are rational, then \(L_x\) and \(M_x\) are rational, the guarded
ratio defining \(s_x\) is rational, and the entire attaining model is
rational.

The product in (11) is a singleton when \(\lvert\mathcal Y\rvert=1\).
Suppose instead that \(\lvert\mathcal Y\rvert\geq2\).  If every
\(M_x=0\), each factor equals \(\{L_x\}\).  Conversely, if
\(M_{x_0}>0\), two distinct point masses in \(\Delta(\mathcal Y)\),
inserted as \(s_{x_0}\) in (13), give two compatible models with the same
observed law and different \(Q_M(x_0,\cdot)\).  Thus point identification
holds exactly when every \(M_x=0\).  Since every summand defining
\(M_x\) is strictly positive, this is equivalent to
\(p(z,x)>0\) for all \(z\) with \(p(z)>0\) and all \(x\).

Finally, (12) implies the exact equivalence
\[
 q_x\in L_x+M_x\Delta(\mathcal Y)
 \quad\Longleftrightarrow\quad
 q_x(y)\geq L_x(y)\ \text{for all }y,
 \quad \sum_yq_x(y)=1.
\tag{15}
\]
Indeed, the forward implication is immediate.  In the reverse direction,
\(q_x-L_x\) is coordinatewise nonnegative and has coordinate sum \(M_x\);
divide by \(M_x\) when it is positive, while for \(M_x=0\) nonnegativity
and zero sum force \(q_x=L_x\).  This proves the stated linear-size
half-space representation.  Computing all marginal cells requires one pass
through \(\lvert\mathcal Z\rvert\lvert\mathcal X\rvert
\lvert\mathcal Y\rvert\) cells.  Subject to their explicit strict guards,
the ratios, profile entries, \(L_x(y)\), and \(M_x\) require the same order
of arithmetic operations.  After exact profile comparison, counting
distinct profiles and testing collapse are linear in the stored profile
table.  The representation (15) has
\(\lvert\mathcal X\rvert\lvert\mathcal Y\rvert\) inequalities and
\(\lvert\mathcal X\rvert\) equalities, completing the claimed bounds.